\section{User Guide} \label{sec:user-guide}

The project can be executed either with Docker Compose or manually by starting the backend and frontend separately.
The recommended method is Docker Compose, because it starts both services with a single command and avoids dependency conflicts.

\subsection{Docker Setup} \label{subsec:docker-setup}

The repository includes a \texttt{docker-compose.yml} file at the project root.
It defines two services:

\begin{itemize}
    \item \texttt{backend}: builds the FastAPI backend from \texttt{./backend} and exposes it on port \texttt{8000}.
    \item \texttt{frontend}: builds the Vue/Vite frontend from \texttt{./frontend} and exposes it on port \texttt{3000}.
\end{itemize}

\noindent To start the complete application using Docker, run the following command from the root directory of the project:

\begin{lstlisting}[language=bash,label={lst:start-docker}]
docker-compose up --build
\end{lstlisting}

\noindent After the containers are running, the application can be accessed at:

\begin{lstlisting}[language=bash, label={lst:endpoint-link}]
Frontend UI:      http://localhost:3000
Backend API docs: http://localhost:8000/docs
Backend API:      http://localhost:8000/api
\end{lstlisting}

\noindent To stop the application, press:

\begin{lstlisting}[language=bash, label={lst:stop-application}]
Ctrl+C
\end{lstlisting}

\noindent Alternatively, to stop and remove the containers, run:

\begin{lstlisting}[language=bash,label={lst:stop-docker}]
docker-compose down
\end{lstlisting}

\noindent The frontend container is configured with the environment variable:

\begin{lstlisting}[language=bash,label={lst:frontend-env}]
VITE_API_BASE_URL=http://localhost:8000/api
\end{lstlisting}

\noindent This means that, when the frontend is opened in the browser at \texttt{localhost:3000}, API requests are sent to the backend at \texttt{localhost:8000/api}.

\subsection{Backend Setup Without Docker} \label{subsec:backend-setup-without-docker}

To run the backend manually:

\begin{lstlisting}[language=bash,label={lst:run-backend}]
cd backend
python -m venv .venv
source .venv/bin/activate
pip install -r requirements.txt
uvicorn app.main:app --reload
\end{lstlisting}

\noindent On Windows, the virtual environment can be activated with:

\begin{lstlisting}[language=bash,label={lst:create-venv-window}]
.venv\Scripts\activate
\end{lstlisting}

\noindent The backend runs at:

\begin{lstlisting}[language=bash,label={lst:backend-link}]
http://localhost:8000
\end{lstlisting}

\noindent The API documentation is available at:

\begin{lstlisting}[language=bash,label={lst:api-doc-link}]
http://localhost:8000/docs
\end{lstlisting}

\subsection{Frontend Setup Without Docker} \label{subsec:frontend-setup-without-docker}

To run the frontend manually:

\begin{lstlisting}[language=bash, label={lst:run-frontend}]
cd frontend
npm install
npm run dev
\end{lstlisting}

\noindent When running through Docker, the frontend is available at:

\begin{lstlisting}[language=bash, label={lst:link-frontend-docker}]
http://localhost:3000
\end{lstlisting}

\subsection{CLI Examples} \label{subsec:cli-examples}

Task 1 example:

\begin{lstlisting}[language=bash,label={lst:example-1}]
python cli.py --floors 3 --persons 3 --task 1 --goal-floor 2
\end{lstlisting}

\noindent Task 2 example:

\begin{lstlisting}[language=bash,label={lst:example-2}]
python cli.py --floors 3 --persons 2 --items chair:1 table:2
\end{lstlisting}

\noindent Fixed-horizon solving:

\begin{lstlisting}[language=bash,label={lst:example-3}]
python cli.py --floors 3 --persons 2 --items chair:1 table:2 --fixed-t 5
\end{lstlisting}

\noindent Exporting DIMACS CNF:

\begin{lstlisting}[language=bash,label={lst:exporting-dimacs}]
python cli.py --floors 3 --persons 2 --items chair:1 table:2 --cnf-out problem.cnf
\end{lstlisting}